\documentclass[11pt,a4paper]{article}
\usepackage[utf8]{inputenc}
\usepackage[T1]{fontenc}
\usepackage[portuguese]{babel}
\usepackage{geometry}
\geometry{margin=2.5cm}
\usepackage{listings}
\usepackage{xcolor}
\usepackage{amsmath}
\usepackage{amssymb}
\usepackage{hyperref}
\usepackage{enumitem}
\usepackage{tcolorbox}
\usepackage{tabularx}
\usepackage{booktabs}
\usepackage{graphicx}

\definecolor{codebg}{rgb}{0.95,0.95,0.95}
\definecolor{codegreen}{rgb}{0.1,0.5,0.1}
\definecolor{codepurple}{rgb}{0.5,0.1,0.5}
\definecolor{codegray}{rgb}{0.5,0.5,0.5}

\lstdefinestyle{python}{
    backgroundcolor=\color{codebg},
    commentstyle=\color{codegray}\itshape,
    keywordstyle=\color{codepurple}\bfseries,
    stringstyle=\color{codegreen},
    basicstyle=\ttfamily\small,
    breaklines=true,
    frame=single,
    language=Python,
    showstringspaces=false,
    tabsize=4,
    literate={á}{{\'a}}1 {ã}{{\~a}}1 {é}{{\'e}}1 {í}{{\'i}}1 {ó}{{\'o}}1 {õ}{{\~o}}1 {ú}{{\'u}}1 {ç}{{\c{c}}}1 {Á}{{\'A}}1 {Ã}{{\~A}}1 {É}{{\'E}}1 {Í}{{\'I}}1 {Ó}{{\'O}}1 {Õ}{{\~O}}1 {Ú}{{\'U}}1 {Ç}{{\c{C}}}1 {â}{{\^a}}1 {ê}{{\^e}}1 {ô}{{\^o}}1 {û}{{\^u}}1 {Â}{{\^A}}1 {Ê}{{\^E}}1 {Ô}{{\^O}}1 {Û}{{\^U}}1 {à}{{\`a}}1 {è}{{\`e}}1 {ì}{{\`i}}1 {ò}{{\`o}}1 {ù}{{\`u}}1 {ä}{{\"a}}1 {ö}{{\"o}}1 {Ä}{{\"A}}1 {Ö}{{\"O}}1
}
\lstset{style=python}

\newtcolorbox{objetivos}[1][]{
  colback=green!5!white,
  colframe=green!40!black,
  title=O que você vai aprender nesta página,
  fonttitle=\bfseries,
  #1
}

\title{\textbf{Data Analysis with Python 2024--2025}\\Parte 7: Considerações Finais\\\large Conclusão do Curso e Próximos Passos}
\author{Tradução e Adaptação: University of Helsinki --- MOOC.fi}
\date{}

\begin{document}

\maketitle

\section*{Nota Importante}
Este material é uma tradução e adaptação baseada no curso \textit{Data Analysis with Python 2024-2025}, oferecido pela \textbf{Universidade de Helsinque (MOOC.fi)}.

\tableofcontents
\newpage

\section{Considerações Finais (\emph{Closing Words})}

Neste curso, exploramos diferentes maneiras de analisar dados utilizando a linguagem Python. Utilizamos diversas bibliotecas para estruturar, visualizar e realizar avaliações estatísticas em variados conjuntos de dados. Também aprendemos modelos úteis de aprendizado de máquina (\emph{ML}) que são aplicados para aprender a partir de dados e realizar previsões ou agrupamentos (\emph{clustering}). 

No entanto, esta é apenas a base de uma longa jornada rumo ao domínio completo das habilidades de análise de dados.

\section{Próximos Passos e Recomendações}

Existem muitos cursos bons de Python disponíveis online que qualquer pessoa pode utilizar para fortalecer suas competências de programação. Para a parte específica de ciência de dados, a Universidade de Helsinque recomenda os seguintes cursos (atualmente disponíveis apenas para alunos regulares da instituição):

\begin{itemize}
    \item \textbf{Introduction to Data Science} (Introdução à Ciência de Dados);
    \item \textbf{Introduction to Machine Learning} (Introdução ao Aprendizado de Máquina);
    \item \textbf{Bayesian Data Analysis} (Análise de Dados Bayesianos).
\end{itemize}

\section{Conclusão do Curso}

Para estar apto a realizar a prova final e submeter o projeto de extensão, o estudante deve ter:
\begin{enumerate}
    \item Estudado todo o material do curso;
    \item Concluído os exercícios relacionados.
\end{enumerate}

Para finalizar o curso com êxito, é necessário:
\begin{enumerate}
    \item Concluir a prova (\emph{exam});
    \item Submeter o projeto prático;
    \item Enviar e receber o feedback dos pares (\emph{peer feedback}).
\end{enumerate}

Parabéns por concluir o material didático deste curso!

\end{document}
